% !TEX root = frode-bibelen.tex
% ----------------------------------------------------------------------------
%  Frode-bibelen -- forside og kolofon
%  Inkluderes fra frode-bibelen.tex (kompiler alltid hovedfilen).
% ----------------------------------------------------------------------------

% =============================================================== FORSIDE ===
\hypersetup{pageanchor=false}
\begin{titlepage}
\thispagestyle{empty}
\ifmorkforside
  \colorlet{omslag}{skinn}\colorlet{omslaglys}{skinnlys}
  \colorlet{forsidetekst}{gull}\colorlet{ramme}{gull}\colorlet{ramme2}{gullmork}
\else
  \colorlet{omslag}{white}\colorlet{omslaglys}{white}
  \colorlet{forsidetekst}{skinn}\colorlet{ramme}{gullmork}\colorlet{ramme2}{gull}
\fi
\newsavebox{\asciiboks}\newsavebox{\asciiskygge}
\begin{lrbox}{\asciiboks}\begin{minipage}{7.2cm}
\begin{lstlisting}[language={},backgroundcolor={},frame=none,xleftmargin=0pt,framexleftmargin=0pt,aboveskip=0pt,belowskip=0pt,columns=fixed,keepspaces=true,upquote=true,breaklines=false,literate={},basicstyle=\ttfamily\bfseries\fontsize{11.5}{12.4}\selectfont\color{forsidetekst}]
           cC
          cC
          cC
           cC
            cCc
              cC
               cC
                cC
                cC
                cC
c               cC
Cc     .-'\    cC
 Cc   /\   \  cC
  Cc. \ \   \cC  .-'\
   `Cc.\ \   \.-'    \
     `Cc\ \        .-/
        .-'     .-'.-'
     .-'        \-'
    /\    .-/\   \
    \ \.-'.-\ \   \
     \/.-'   \ \   \
              \ \   \
               \ \   \
                \ \   \
                 \ \   \
                  \ \   \
                   \ \   \
                    \ \_-/
                     \/-'
\end{lstlisting}
\end{minipage}\end{lrbox}
\begin{lrbox}{\asciiskygge}\begin{minipage}{7.2cm}
\begin{lstlisting}[language={},backgroundcolor={},frame=none,xleftmargin=0pt,framexleftmargin=0pt,aboveskip=0pt,belowskip=0pt,columns=fixed,keepspaces=true,upquote=true,breaklines=false,literate={},basicstyle=\ttfamily\bfseries\fontsize{11.5}{12.4}\selectfont\color{black}]
           cC
          cC
          cC
           cC
            cCc
              cC
               cC
                cC
                cC
                cC
c               cC
Cc     .-'\    cC
 Cc   /\   \  cC
  Cc. \ \   \cC  .-'\
   `Cc.\ \   \.-'    \
     `Cc\ \        .-/
        .-'     .-'.-'
     .-'        \-'
    /\    .-/\   \
    \ \.-'.-\ \   \
     \/.-'   \ \   \
              \ \   \
               \ \   \
                \ \   \
                 \ \   \
                  \ \   \
                   \ \   \
                    \ \_-/
                     \/-'
\end{lstlisting}
\end{minipage}\end{lrbox}
\begin{tikzpicture}[remember picture,overlay]
  % Skinnomslag med svak lysning på midten
  \shade[inner color=omslaglys, outer color=omslag]
        (current page.south west) rectangle (current page.north east);
  % Falsen ved ryggen (venstre kant)
  \ifmorkforside
    \fill[black,opacity=.35] ([xshift=0.95cm]current page.south west)
                       rectangle ([xshift=1.25cm]current page.north west);
    \draw[white,opacity=.07,line width=.6pt] ([xshift=1.3cm]current page.south west)
                       -- ([xshift=1.3cm]current page.north west);
  \fi
  % Doble gullrammer
  \coordinate (SW) at ([shift={(2.0cm,1.4cm)}]current page.south west);
  \coordinate (NE) at ([shift={(-1.4cm,-1.4cm)}]current page.north east);
  \coordinate (sw) at ([shift={(2.3cm,1.7cm)}]current page.south west);
  \coordinate (ne) at ([shift={(-1.7cm,-1.7cm)}]current page.north east);
  \draw[ramme,line width=1.8pt] (SW) rectangle (NE);
  \draw[ramme2,line width=.6pt] (sw) rectangle (ne);
  % Hjørneornamenter
  \coordinate (NW) at (SW|-NE);  \coordinate (SE) at (NE|-SW);
  \foreach \p in {SW,NE,NW,SE}
    \fill[ramme,rotate around={45:(\p)}] ($(\p)+(-.2,-.2)$) rectangle ++(.4,.4);
  % Midtpunkt for alt innhold (midt i rammen)
  \coordinate (M) at ($(SW)!.5!(NE)$);
  % Fleuron midt på øvre og nedre ramme
  \node[fill=omslag,inner sep=4pt,text=ramme,font=\Huge] at (M|-NE) {\ding{167}};
  \node[fill=omslag,inner sep=4pt,text=ramme,font=\Huge,rotate=180] at (M|-SW) {\ding{167}};

  % Øverst
  \node[text=forsidetekst,font=\fontfamily{ppl}\selectfont\small] at ($(M|-ne)+(0,-1.4)$)
       {\textls[250]{\MakeUppercase{Prog1001 \,\textperiodcentered\, Grunnleggende programmering}}};

  % ASCII-motivet (preget i gull, med skygge på mørk forside)
  \ifmorkforside
    \node[inner sep=0pt,opacity=.55] at ($(M)+(0.63,3.17)$) {\usebox{\asciiskygge}};
  \fi
  \node[inner sep=0pt] (A) at ($(M)+(0.6,3.2)$) {\usebox{\asciiboks}};

  % Tittel
  \node[text=forsidetekst,font=\fontfamily{ppl}\fontsize{44}{48}\selectfont] (T)
       at ($(M)+(0,-4.5)$) {\textls[120]{FRODE-BIBELEN}};
  % Ornamentlinje
  \draw[ramme,line width=.7pt] ($(T.south)+(-5.2,-.45)$) -- ($(T.south)+(-.35,-.45)$);
  \draw[ramme,line width=.7pt] ($(T.south)+(.35,-.45)$) -- ($(T.south)+(5.2,-.45)$);
  \fill[ramme,rotate around={45:($(T.south)+(0,-.45)$)}] ($(T.south)+(-.12,-.57)$) rectangle ++(.24,.24);
  % Undertittel
  \node[text=forsidetekst,font=\fontfamily{ppl}\selectfont\Large\itshape,align=center]
       at ($(T.south)+(0,-1.35)$) {Den hellige skrift for grunnleggende programmering i C};
  % Nederst
  \node[text=forsidetekst,font=\fontfamily{ppl}\selectfont\normalsize] at ($(M|-sw)+(0,1.9)$)
       {\textls[250]{\MakeUppercase{Eksamensutgaven}}};
  \node[text=forsidetekst,font=\fontfamily{ppl}\selectfont\small\itshape] at ($(M|-sw)+(0,1.25)$)
       {med maler, fallgruver og gjennomregnede eksempler \textperiodcentered\ 2026};
\end{tikzpicture}
\end{titlepage}
\hypersetup{pageanchor=true}

% ============================================================== KOLOFON ===
\thispagestyle{empty}
\vspace*{1.5cm}
\begin{flushright}
{\fontfamily{ppl}\selectfont\Large\itshape «I begynnelsen var \texttt{main()}, og \texttt{main()} var hos Frode.»}\\[4pt]
{\fontfamily{ppl}\selectfont\small -- fritt etter Joh 1,1}
\end{flushright}
\vspace{1.5cm}

\begin{regel}[Slik bruker du bibelen]
\begin{itemize}
  \item \textbf{Skriv den ut} og legg den i en perm med faneark per kapittel. Topptekstene viser hvilket kapittel du er i.
  \item \textbf{Kapittel 8--9} (maler og filer) er det du blar mest i under oppgave 2. \textbf{Kapittel 11} og \textbf{vedleggene} (ASCII og sporingsark) er for oppgave 1.
  \item \textbf{Ta med originalen av \cc{LesData.h}} -- Frode anbefaler selv en utskrift. Legger du filen i samme mappe som \cc{.tex}-filen, kommer den automatisk med i kapittel 2.
  \item \textbf{Ta med 2--3 løsningsforslag} (\cc{ex\_xxx\_2.c}) fra eksamensarkivet. De viser hele program i nøyaktig den stilen sensor kjenner igjen.
  \item \textbf{Kodemalene her er egne maler i Frodes stil} med \emph{Kunde} som eksempel. Bytt ut \cc{Kunde}/\cc{kunde} med det oppgaven handler om.
  \item \textbf{Fyll på selv} i margen etter hver øvingseksamen.
\end{itemize}
\end{regel}

\vfill
{\small Bygget på Frode Haugs emneside
\href{https://frh.folk.ntnu.no/grprog/}{\cc{frh.folk.ntnu.no/grprog}}:
opplegg, eksamensråd, stikkordsoppsummering, eksempler (EKS\_01--EKS\_39),
\cc{LesData.h} og eksamenene H19--S26. Uoffisielt studentlaget hjelpemiddel --
ikke laget eller godkjent av emneansvarlig.\par}
\vspace{0.3cm}
{\small\sffamily Eksamensutgaven, november 2026}\\[2pt]
{\small\sffamily Kildekode og lisens (Beerware, Revision 42):
 \href{https://github.com/mikkeliski/frode-bibelen}{\cc{github.com/mikkeliski/frode-bibelen}}}
\clearpage
